%
% Das Integrationsproblem
%
\subsection{Das Integrationsproblem
\label{buch:intalg:wurzel:subsection:integration}}
Die in Abschnitt~\ref{buch:intalg:wurzel:subsection:rationalisierung}
erfolgte Rationalisierung zeigt, dass eine rationale Funktion der Form
$R(x,\vartheta)$ immer in die Form \eqref{buch:intalg:wurzel:eqn:linkombtheta}
oder \eqref{buch:intalg:wurzel:eqn:linkombthetainverse} gebracht
werden kann.
Es müssen also nur noch Integrale der Form
\begin{equation*}
\int
p(x)+q(x)\vartheta
\,dx
\qquad\text{oder}\qquad
\int
p(x)+q(x)\frac{1}{\vartheta}
\,dx
\end{equation*}
gelöst werden, wobei $p$ und $q$ rationale Funktionen in $\mathbb{Q}(x)$
sind.
Das Integral über den ersten Summanden ist das Integral einer rationalen
Funktion und kann daher mit der Methode von Kapitel~\ref{chapter:ratint}
integriert werden.
Es bleibt daher nur noch Integrale der Form
\begin{equation}
\int q(x)\,\vartheta\,dx
\qquad\text{oder}\qquad
\int \frac{q(x)}{\vartheta}\,dx
\label{buch:intalg:wurzel:eqn:integrand}
\end{equation}
zu bestimmen.


Im Folgenden bevorzugen wir die zweite Form des Integrals in
\eqref{buch:intalg:wurzel:eqn:integrand}.
Um den Grund zu verstehen, betrachten wir die Ableitung
\begin{align*}
\vartheta^2
&=
ax^2+bx+c
\\
2\vartheta\vartheta'
&=
2ax + b
\intertext{und drücken $\vartheta'$ durch $\vartheta$ als}
\vartheta'
&=
\frac{2ax+b}{2\vartheta}
=
\frac{ax+b/2}{\vartheta}
&&\Rightarrow&
\frac{1}{\vartheta}
&=
\frac{\vartheta'}{ax+b/2}
\end{align*}
aus.
Man kann also $1/\vartheta$ durch die Ableitung $\vartheta'$ ersetzen,
was die Möglichkeit der partiellen Integration eröffnet.

Die rationale Funktion $q(x)$ im rechten Integral von
\eqref{buch:intalg:wurzel:eqn:integrand} ist ein Bruch von Polynomen
von $x$.
Durch Dividieren mit Rest lässt es sich als Summe
\begin{equation}
q(x)
=
\varphi(x)
+
w(x)
\qquad
\varphi(x)\in\mathbb{Q}[x],
w(x)\in\mathbb{Q}(x)
\label{buch:intalg:wurzel:eqn:summe}
\end{equation}
eines Polynoms und einer vollständig gekürzten rationalen Funktion
schreiben, deren Nenner grösseren Grad hat als der Zähler.
